% Part: proof-theory
% Chapter: cut-elimination
% Section: ce-topmost

\documentclass[../../../include/open-logic-section]{subfiles}

\begin{document}

\olfileid{pt}{cut}{top}

\olsection{Ngilangi Cut kang Paling Dhuwur}

\begin{defn}
Gatekna !!a{proof}~$\pi$ kang pungkasané nganggo aturan \CutCS{},
\[
\AxiomC{}
\RightLabel{$\pi_1$}
\Deduce$\Gamma \fCenter \Delta, !A$
\AxiomC{}
\RightLabel{$\pi_2$}
\Deduce$!A, \Gamma \fCenter \Delta$
\RightLabel{\CutCS}
\BinaryInf$\Gamma \fCenter \Delta$
\DisplayProof
\]
lan ora ngemot inferènsi \CutCS{} liyané. \emph{Dhuwuré cut}
saka~$\pi$ yaiku $\cheight{\pi} = \pheight{\pi_1} + \pheight{\pi_2}$.
\emph{Pangkat cut} saka~$\pi$ yaiku $\cutrank{\pi} = \depth{!A}$.
\end{defn}

\begin{lem}\ollabel{lem:cut-adm-G3c}
Aturan \CutCS{} admisibel ing~\Log{G3c}. Tegesé, yèn $\pi$
iku !!a{proof} kang pungkasané nganggo siji~\CutCS{} lan saliyané
mung nganggo aturan~$\Log{G3c}$, ana !!a{proof}~$\pi'$ ing
\Log{G3c} kanggo sekèn pungkasan kang padha.
\end{lem}

\begin{proof}
Buktiné nganggo induksi rangkep tumrap !!{height} lan pangkat cut.
!!{proof}~$\pi$ pungkasané kaya mangkéné:
\[
\AxiomC{}
\RightLabel{$\pi_1$}
\Deduce$\Gamma \fCenter \Delta, !A$
\AxiomC{}
\RightLabel{$\pi_2$}
\Deduce$!A, \Gamma \fCenter \Delta$
\RightLabel{\CutCS}
\BinaryInf$\Gamma \fCenter \Delta$
\DisplayProof
\]

Dhasar induksi: $\cheight{\pi} = 0$ lan $\cutrank{\pi} = 0$.
Amarga $\cheight{\pi} = \pheight{\pi_1} + \pheight{\pi_2} = 0$,
$\Gamma \Sequent \Delta, !A$ lan $!A, \Gamma \Sequent \Delta$
kaloroné aksioma. Ana rong kahanan: (a)~$!A$ dudu formula utama
ing salah siji aksioma, utawa (b)~$!A$ formula utama ing kaloroné.
Ing saben kahanan, kita bakal nuduhaké ing ngisor yèn
$\Gamma \Sequent \Delta$ mesthi aksioma (tanpa migunakaké
hipotesis induksi).

Saiki kita bedakaké sawatara kahanan manut $\pi_1$ lan
$\pi_2$ iku aksioma utawa pungkasané nganggo inferènsi, lan
manut $!A$~iku formula utama ing inferènsi mau utawa ora.
Kahanan A lan~B nyakup dhasar induksi. Ing langkah induksi,
kita anggep $\cheight{\pi} > 0$ utawa $\cutrank{\pi} > 0$
(utawa kaloroné). Kita bisa migunakaké hipotesis induksi:
\CutCS{} pungkasan bisa diilangi saka saben !!{proof} kang
pungkasané nganggo siji~\CutCS{} lan pangkat cuté
$= \cutrank{\pi}$ kanthi !!{height} cut $< \cheight{\pi}$,
utawa pangkat cuté~$< \cutrank{\pi}$. Kahanan
$\cheight{\pi} = 0$ (kaloro premis aksioma) katangani
dening kahanan~A lan~B. Kahanan $\cheight{\pi} > 0$
katangani dening kahanan C--D.

\paragraph{Kahanan A:}
Ana $!B$ ing $\Gamma$ lan uga ing~$\Delta$.
Mula $\Gamma \Sequent \Delta$ iku aksioma (yèn $!B$~atomik),
utawa nduwèni !!a{proof}~$\pi'$ ing \Log{G3c} tanpa
\CutCS{} miturut \olref[seq][ex]{prop:G3c-axioms}.
Iki nyakup kahanan~(a) ing dhasar induksi: yèn $!A$ dudu
formula utama ing $\Gamma \Sequent \Delta, !A$ utawa ing
$!A, \Gamma \Sequent \Delta$, lan kaloroné aksioma, mesthi
ana $!B$~atomik liyané ing $\Gamma$ lan~$\Delta$.
Iki uga nyakup langkah induksi yèn salah siji premis
iku aksioma kang $!A$ dudu formula utamané (sanadyan
premis liyané minangka konklusi sawijining aturan).

\paragraph{Kahanan B:} Kaloro premis iku aksioma, lan $!A$
formula utama, mula atomik. Gatekna premis kiwa,
$\Gamma \Sequent \Delta, !A$. Amarga $!A$ formula utama,
formula iku kudu ana ing $\Gamma$ (yèn ora,
$\Gamma \Sequent \Delta, !A$ dudu aksioma). Semono uga,
$!A$ kudu ana ing $\Delta$ (yèn ora,
$!A, \Gamma \Sequent \Delta$ dudu aksioma).
Dadi $!A$ atomik lan ana ing $\Gamma$ uga $\Delta$;
tegesé, $\Gamma \Sequent \Delta$ iku aksioma. Iki
nyakup kahanan~(b) ing dhasar induksi.

\begin{center}
Kahanan A lan B versi alternatif
\end{center}

\paragraph{Kahanan A:} Salah siji premis cut iku aksioma
kanthi wujud $!C, \Gamma' \Sequent \Delta', !C$ lan $!C$~atomik.
Yèn $!C$ dudu !!{formula} cut~$!A$, banjur $!C$ mesthi
ana ing $\Gamma$ lan~$\Delta$. Mula $\Gamma \Sequent \Delta$
iku aksioma lan bisa dijupuk minangka~$\pi'$.
Yèn $!C \ident !A$ iku !!{formula} cut, banjur
(yèn premis kiwa aksioma) $!A \in \Gamma$ lan
$\Gamma = \Gamma', !A$, utawa (yèn premis tengen aksioma)
$!A \in \Delta$ lan $\Delta = \Delta', !A$.
Ing kahanan kapisan, $\pi_2$ iku !!a{proof} kanggo
$!A, !A, \Gamma' \Sequent \Delta$ lan miturut admisibilitas
\LeftR{\Contraction} (\olref[seq][inv]{prop:G3c-cont-adm})
ana !!a{proof} kanggo $\Gamma \Sequent \Delta$.
Ing kahanan kapindho, $\pi_1$ iku !!a{proof} kanggo
$\Gamma \Sequent \Delta', !A, !A$ lan miturut admisibilitas
\RightR{\Contraction} kita ngolèhké !!a{proof} kanggo
$\Gamma \Sequent \Delta$.

\paragraph{Kahanan B:} Salah siji premis iku aksioma kanthi
wujud~$\lfalse, \Gamma' \Sequent \Delta'$. Yèn premis kiwa
aksioma mangkono, banjur $\lfalse \in \Gamma$ lan
$\Gamma \Sequent \Delta$ uga aksioma. Yèn premis tengen
aksioma mangkono, banjur $\lfalse \in \Gamma$
(mula $\Gamma \Sequent \Delta$ uga aksioma), utawa
$!A \ident \lfalse$. Ing kahanan kapindho iki, $\pi_1$
iku !!a{proof} kanggo
$\Gamma \Sequent \Delta, \lfalse$. Bukti iku bisa diowahi
dadi !!a{proof} kanggo~$\Gamma \Sequent \Delta$ kanthi
mbusak siji salinan~$\lfalse$ saka suksèden saben sekèn
ing~$\pi_1$. (Latihan: priksanen kanthi induksi tumrap
$\pheight{\pi_1}$.)

Saiki anggepen kahanan~A utawa~B ora lumaku. Ing kahanan
liyané, $\cheight{\pi} > 0$, yaiku paling ora salah siji
premis minangka konklusi sawijining inferènsi.

\paragraph{Kahanan C lan D: Ngowahi urutan cut.}
Kabèh kahanan ing~C lan~D nduwèni pola kang padha.
Kita miwiti saka !!a{proof} kang pungkasané nganggo~\CutCS{},
kanthi salah siji premis diolehké saka aturan~$R$ kang
!!{formula} cut~$!A$ dudu formula utamané (nanging
!!{formula} liya~$!D$). Ing !!{proof} iki,
\CutCS{} ana sawisé aturan~$R$, kanthi skéma mangkéné:
\[
\AxiomC{}
\RightLabel{$\pi_1$}
\Deduce$\Gamma \fCenter \Delta', !D, !A$
\AxiomC{}
\RightLabel{$\pi_3$}
\Deduce$!A, \Gamma, \Pi \fCenter \Delta', \Lambda$
\RightLabel{$R$}
\UnaryInf$!A, \Gamma \fCenter \Delta', !D$
\RightSubproofLabel{$\pi_2$}
\RightLabel{\CutCS}
\BinaryInf$\Gamma \fCenter \Delta', !D$
\DisplayProof
\]
Iki mung nyakup kahanan nalika $R$~aturan tengen kang nduwèni
siji premis, $A$~dudu !!{formula} utama saka~$R$, lan
!!{formula}~$!D$ kang \emph{dadi} formula utama ana ing
suksèden premis tengen~\CutCS. Yèn $!D$~ana ing antesèden
(yaiku $R$ aturan kiwa) utawa ing premis kiwa~\CutCS,
wujudé mesthi béda. !!{formula}s ing $\Pi$ lan~$\Lambda$
makili !!{formula}s aktif saka aturan~$R$.

Kita balik urutan iki lan gatekna !!a{proof}
ing ngendi $R$ ana sawisé~\CutCS{}:
\[
\AxiomC{}
\RightLabel{$\pi_1'$}
\Deduce$\Gamma, \Pi \fCenter \Delta', !D, \Lambda, !A$
\AxiomC{}
\RightLabel{$\pi_3'$}
\Deduce$!A, \Gamma, \Pi \fCenter \Delta', !D, \Lambda$
\RightLabel{\CutCS}
\BinaryInf$\Gamma, \Pi \fCenter \Delta', !D, \Lambda$
\RightLabel{$R$}
\UnaryInf$\Gamma \fCenter \Delta', !D, !D$
\DisplayProof
\]
Kita migunakaké pelemahan kang njaga !!{height} kanggo ngolèhké
$\pi_1'$ lan $\pi_3'$ saka $\pi_1$ lan $\pi_3$ miturut urutané.

Amarga urutan \CutCS{} lan~$R$ dibalik, iki diarani
``ngowahi urutan cut munggah ngliwati~$R$.'' Tumindak iki
nyuda dhuwuré !!{proof} kang pungkasané ana ing premis
tengen~\CutCS, mula uga nyuda !!{height} cut. Tegesé,
sub-!!{proof}s kang pungkasané ana ing \CutCS{} anyar bisa
dianggep !!{height}-é ora ngluwihi $\pi_1$ lan $\pi_3$,
mula !!{height} cut suda (siji inferènsi ing sisih tengen
wis kabusak). Saiki hipotesis induksi bisa ditrapaké,
mula inferènsi \CutCS{} bisa diilangi. Pungkasané,
salinan ekstra saka~$!D$ dibusak nganggo admisibilitas
$\RightR{\Contraction}$.

\paragraph{Kahanan C:}
Paling ora salah siji saka $\pi_1$ lan $\pi_2$ pungkasané
nganggo aturan siji premis~$R$ kang $!A$~dudu formula utama.
Gunggungé ana 18 kahanan, manut $!A$~dudu formula utama
ing premis kiwa utawa tengen~\CutCS, sarta manut endi
saka sangang aturan siji premis
(\LeftR{\lnot}, \RightR{\lnot}, \RightR{\lor}, \LeftR{\land},
\RightR{\lif}, lan papat aturan kuantor) kang dadi~$R$.
Kita mung ngrembug rong tuladha: $!A$~dudu formula utama
ing inferènsi pungkasan~$\pi_2$ kang nganggo~$\RightR{\lif}$
utawa~\LeftR{\lexists}. Kahanan liyané padha polané
lan dipasrahaké dadi latihan.

Anggepen $\pi$ pungkasané kaya mangkéné:
\[
\AxiomC{}
\RightLabel{$\pi_1$}
\Deduce$\Gamma \fCenter \Delta', !B \lif !C, !A$
\AxiomC{}
\RightLabel{$\pi_3$}
\Deduce$!A, !B, \Gamma \fCenter \Delta', !C$
\RightLabel{$\RightR{\lif}$}
\UnaryInf$!A, \Gamma \fCenter \Delta', !B \lif !C$
\RightSubproofLabel{$\pi_2$}
\RightLabel{\CutCS}
\BinaryInf$\Gamma \fCenter \Delta', !B \lif !C$
\DisplayProof
\]
ing ngendi $\Delta = \Delta', !B \lif !C$. Kita ngetrapaké
pelemahan kang njaga !!{height}
(\olref[seq][adm]{prop:weak-G3c-adm}) marang $\pi_1$
kanggo ngolèhké !!a{proof} $\pi_1'$ kanggo
$!B, \Gamma \fCenter \Delta', !B \lif !C, !C, !A$
(kanthi nambah $!B$ ing kiwa lan $!C$ ing tengen).
Semono uga, kita ngetrapaké
\olref[seq][adm]{prop:weak-G3c-adm} marang~$\pi_3$
kanggo ngolèhké !!a{proof}~$\pi_3'$ kanggo
$!A, !B, \Gamma \fCenter \Delta', !B \lif !C, !C$
(kanthi nambah $!B \lif !C$ ing tengen).
Hipotesis induksi ditrapaké marang
\[
\AxiomC{}
\RightLabel{$\pi_1'$}
\Deduce$!B, \Gamma \fCenter \Delta', !B \lif !C, !C, !A$
\AxiomC{}
\RightLabel{$\pi_3'$}
\Deduce$!A, !B, \Gamma \fCenter \Delta', !B \lif !C, !C$
\RightLabel{\CutCS}
\BinaryInf$!B, \Gamma \fCenter \Delta', !B \lif !C, !C$
\DisplayProof
\]
kanthi !!{formula} cut~$!A$: !!{proof} iki nduwèni
pangkat cut kang padha nanging !!{height} cut luwih cilik
tinimbang~$\pi$. Mula kita ngolèhké !!a{proof}~$\pi_4$
ing \Log{G3c} tanpa \CutCS{} kanggo
$!B, \Gamma \fCenter \Delta', !B \lif !C, !C$.
Trapna $\RightR{\lif}$ kanggo ngolèhké !!a{proof} kanggo
$\Gamma \fCenter \Delta', !B \lif !C, !B \lif !C$:
\[
\AxiomC{}
\RightLabel{$\pi_4$}
\Deduce$!B, \Gamma \fCenter \Delta', !B \lif !C, !C$
\RightLabel{$\RightR{\lif}$}
\UnaryInf$\Gamma \fCenter \Delta', !B \lif !C, !B \lif !C$
\DisplayProof
\]
Kanggo ngolèhké !!a{proof}~$\pi'$ kanggo
$\Gamma \Sequent \Delta$, yaiku
$\Gamma \fCenter \Delta', !B \lif !C$, trapna
\olref[seq][inv]{prop:G3c-cont-adm} (admisibilitas kontraksi).

Saiki anggepen $\pi$ pungkasané kaya mangkéné:
\[
\AxiomC{}
\RightLabel{$\pi_1$}
\Deduce$\lexists[x][!B(x)], \Gamma' \fCenter \Delta, !A$
\AxiomC{}
\RightLabel{$\pi_3$}
\Deduce$!A, !B(c), \Gamma' \fCenter \Delta$
\RightLabel{$\LeftR{\lexists}$}
\UnaryInf$!A, \lexists[x][!B(x)], \Gamma' \fCenter \Delta$
\RightSubproofLabel{$\pi_2$}
\RightLabel{\CutCS}
\BinaryInf$\Gamma' \fCenter \Delta$
\DisplayProof
\]
Hipotesis induksi kita trapaké manèh marang
\[
\AxiomC{}
\RightLabel{$\pi_1[!B(c) \Sequent {}]$}
\Deduce$\lexists[x][!B(x)], !B(c), \Gamma' \fCenter \Delta, !A$
\AxiomC{}
\LeftLabel{$\pi_3[\lexists[x][!B(x)] \Sequent {}]$}
\Deduce$!A, \lexists[x][!B(x)], !B(c), \Gamma' \fCenter \Delta$
\RightLabel{\CutCS}
\BinaryInf$\lexists[x][!B(x)], !B(c), \Gamma' \fCenter \Delta$
\RightLabel{\LeftR{\lexists}}
\UnaryInf$\lexists[x][!B(x)], \lexists[x][!B(x)], \Gamma' \fCenter \Delta$
\DisplayProof
\]
Gatekna yèn syarat variabel eigen kacukupan: syarat
variabel eigen kanggo~\LeftR{\lexists} ing~$\pi_2$
njamin $c$ ora ana ing $\lexists[x][!B(x)]$, $\Gamma'$,
utawa~$\Delta$. Pungkasané, trapna
\olref[seq][inv]{prop:G3c-cont-adm}.


\paragraph{Kahanan D:}
Paling ora salah siji saka $\pi_1$ lan $\pi_2$ pungkasané
nganggo aturan rong premis~$R$ kang $!A$ dudu formula utama.
Ana 6 kahanan. Kita ngrembug kahanan nalika $!A$ dudu
formula utama ing inferènsi pungkasan~$\pi_2$ kang nganggo
$\RightR{\land}$; kahanan liyané padha polané.

Anggepen $\pi$ pungkasané kaya mangkéné:
\[
\AxiomC{}
\RightLabel{$\pi_1$}
\Deduce$\Gamma \fCenter \Delta', !B \land !C, !A$
\AxiomC{}
\RightLabel{$\pi_3$}
\Deduce$!A, \Gamma \fCenter \Delta', !B$
\AxiomC{}
\RightLabel{$\pi_4$}
\Deduce$!A, \Gamma \fCenter \Delta', !C$
\RightLabel{$\RightR{\land}$}
\BinaryInf$!A, \Gamma \fCenter \Delta', !B \land !C$
\RightSubproofLabel{$\pi_2$}
\RightLabel{\CutCS}
\BinaryInf$\Gamma \fCenter \Delta$
\DisplayProof
\]
Hipotesis induksi ditrapaké marang !!{proof}s
\[
\AxiomC{}
\RightLabel{$\pi_1'$}
\Deduce$\Gamma \fCenter \Delta', !B \land !C, !B, !A$
\AxiomC{}
\RightLabel{$\pi_3'$}
\Deduce$!A, \Gamma \fCenter \Delta', !B \land !C, !B$
\RightLabel{\CutCS}
\BinaryInf$\Gamma \fCenter \Delta', !B \land !C, !B$
\DisplayProof
\]
and
\[
\AxiomC{}
\RightLabel{$\pi_1''$}
\Deduce$\Gamma \fCenter \Delta', !B \land !C, !C, !A$
\AxiomC{}
\RightLabel{$\pi_4'$}
\Deduce$!A, \Gamma \fCenter \Delta', !B \land !C, !C$
\RightLabel{\CutCS}
\BinaryInf$\Gamma \fCenter \Delta', !B \land !C, !C$
\DisplayProof
\]
Ing kéné, !!{proof}s \Log{G3c} $\pi_1'$ lan $\pi_1''$
diolehké kanthi pelemahan kang njaga !!{height}
(\olref[seq][adm]{prop:weak-G3c-adm}) saka~$\pi_1$,
siji kanthi nambah $!B$ ing tengen, siji manèh kanthi
nambah $!C$. !!{proof}s $\pi_3'$ lan~$\pi_4'$
uga diolehké saka $\pi_3$ lan~$\pi_4$ kanthi
pelemahan kang njaga !!{height} saka $\pi_3$ lan~$\pi_4$
miturut urutané, yaiku nambah
$!B \land !C$ ing tengen.

Hipotesis induksi maringi !!{proof}s~$\pi_5$ lan~$\pi_6$
ing \Log{G3c} kanggo
$\Gamma \fCenter \Delta', !B \land !C, !B$ lan
$\Gamma \fCenter \Delta', !B \land !C, !C$.
Kaloroné digandhèng dadi !!a{proof} kanggo
$\Gamma \Sequent \Delta', !B \land !C, !B \land !C$
nganggo aturan~$\RightR{\land}$:
\[
\AxiomC{}
\RightLabel{$\pi_5$}
\Deduce$\Gamma \fCenter \Delta', !B \land !C, !B$
\AxiomC{}
\RightLabel{$\pi_6$}
\Deduce$\Gamma \fCenter \Delta', !B \land !C, !C$
\RightLabel{\RightR{\land}}
\BinaryInf$\Gamma \fCenter \Delta', !B \land !C, !B \land !C$
\DisplayProof
\]
Kanggo ngolèhké !!a{proof} kanggo
$\Gamma \Sequent \Delta$, yaiku
$\Gamma \fCenter \Delta', !B \land !C$, trapna
\olref[seq][inv]{prop:G3c-cont-adm} (admisibilitas kontraksi).

\paragraph{Kahanan E: Nyuda cut.}
Kahanan pungkasan iki nalika $\pi_1$ lan $\pi_2$ kaloroné
pungkasané nganggo inferènsi kang $!A$ dadi formula utama.
Ana enem kahanan, siji kanggo saben !!{main operator}
kang bisa ana ing~$!A$.

Ing saben kahanan, \CutCS{} kanthi !!{formula} cut~$!A$
diowahi dadi siji utawa luwih cut tumrap sub-!!{formula}s
saka~$!A$, yaiku !!{formula}s aktif saka inferènsi-inferènsi
kang~$!A$ dadi !!{formula} utamané. Mula kahanan iki
biasané diarani ``reduksi'' utawa ``konversi'' inferènsi cut.

\begin{enumerate}
\item Yèn $!A \ident \lnot !B$, banjur $\pi$ pungkasané kaya mangkéné:
\[
\AxiomC{}
\RightLabel{$\pi_1'$}
\Deduce$!B, \Gamma \fCenter \Delta$
\RightLabel{\RightR{\lnot}}
\UnaryInf$\Gamma \fCenter \Delta, \lnot !B$
\LeftSubproofLabel{$\pi_1$}
\AxiomC{}
\RightLabel{$\pi_2'$}
\Deduce$\Gamma \fCenter \Delta, !B$
\RightLabel{\LeftR{\lnot}}
\UnaryInf$\lnot !B, \Gamma \fCenter \Delta$
\RightSubproofLabel{$\pi_2$}
\RightLabel{\CutCS}
\BinaryInf$\Gamma \fCenter \Delta$
\DisplayProof
\]
Miturut hipotesis induksi, \CutCS{} bisa diilangi saka
\[
\AxiomC{}
\RightLabel{$\pi_2'$}
\Deduce$\Gamma \fCenter \Delta, !B$
\AxiomC{}
\RightLabel{$\pi_1'$}
\Deduce$!B, \Gamma \fCenter \Delta$
\RightLabel{\CutCS}
\BinaryInf$\Gamma \fCenter \Delta$
\DisplayProof
\]

\item
Yèn $!A \ident (!B \lor !C)$, banjur $\pi$ pungkasané kaya mangkéné:
\[
\AxiomC{}
\RightLabel{$\pi_1'$}
\Deduce$\Gamma \fCenter \Delta, !B, !C$
\RightLabel{\RightR{\lor}}
\UnaryInf$\Gamma \fCenter \Delta, !B \lor !C$
\LeftSubproofLabel{$\pi_1$}
\AxiomC{}
\RightLabel{$\pi_2'$}
\Deduce$!B, \Gamma \fCenter \Delta$
\AxiomC{}
\RightLabel{$\pi_2''$}
\Deduce$!C, \Gamma \fCenter \Delta$
\RightLabel{\LeftR{\lor}}
\BinaryInf$!B \lor !C, \Gamma \fCenter \Delta$
\RightSubproofLabel{$\pi_2$}
\RightLabel{\CutCS}
\BinaryInf$\Gamma \fCenter \Delta$
\DisplayProof
\]
Gatekna !!{proof} kang pungkasané nganggo \CutCS{},
\[
\AxiomC{}
\RightLabel{$\pi_1'$}
\Deduce$\Gamma \fCenter \Delta, !B, !C$
\AxiomC{}
\RightLabel{$\pi_2'''$}
\Deduce$!C, \Gamma \fCenter \Delta, !B$
\RightLabel{\CutCS}
\BinaryInf$\Gamma \fCenter \Delta, !B$
\DisplayProof
\]
ing ngendi $\pi_2'''$ diolehké saka $\pi_2''$ kanthi
pelemahan kang njaga !!{height}. Pangkat cuté
$< \cutr{\pi}$ amarga $\depth{!C} < \depth{!B \lor !C}$,
mula hipotesis induksi maringi !!a{proof} ing \Log{G3c},
yaiku $\pi_1''$, kanggo $\Gamma \fCenter \Delta, !B$.
!!{proof} kang pungkasané nganggo \CutCS{},
\[
\AxiomC{}
\RightLabel{$\pi_1''$}
\Deduce$\Gamma \fCenter \Delta, !B$
\AxiomC{}
\RightLabel{$\pi_2'$}
\Deduce$!B, \Gamma \fCenter \Delta$
\RightLabel{\CutCS}
\BinaryInf$\Gamma \fCenter \Delta$
\DisplayProof
\]
nduwèni pangkat cut $= \depth{!B} < \depth{!B \lor !C}$,
mula hipotesis induksi bisa ditrapaké manèh lan ngasilaké
!!a{proof} $\pi'$ kanggo $\Gamma \fCenter \Delta$.

\item $!A \ident (!B \land !C)$. Latihan.

\item Yèn $!A \ident (!B \lif !C)$, banjur $\pi$ pungkasané kaya mangkéné:
\[
\AxiomC{}
\RightLabel{$\pi_1'$}
\Deduce$!B, \Gamma \fCenter \Delta, !C$
\RightLabel{\RightR{\lif}}
\UnaryInf$\Gamma \fCenter \Delta, !B \lif !C$
\LeftSubproofLabel{$\pi_1$}
\AxiomC{}
\RightLabel{$\pi_2'$}
\Deduce$\Gamma \fCenter \Delta, !B$
\AxiomC{}
\RightLabel{$\pi_2''$}
\Deduce$!C, \Gamma \fCenter \Delta$
\RightLabel{\LeftR{\lif}}
\BinaryInf$!B \lif !C, \Gamma \fCenter \Delta$
\RightSubproofLabel{$\pi_2$}
\RightLabel{\CutCS}
\BinaryInf$\Gamma \fCenter \Delta$
\DisplayProof
\]
!!{proof} kang pungkasané nganggo \CutCS{},
\[
\AxiomC{}
\RightLabel{$\pi_1'$}
\Deduce$!B, \Gamma \fCenter \Delta, !C$
\AxiomC{}
\RightLabel{$\pi_2''$}
\Deduce$!C, \Gamma \fCenter \Delta$
\RightLabel{\CutCS}
\BinaryInf$!B, \Gamma \fCenter \Delta$
\DisplayProof
\]
nduwèni pangkat cut luwih cilik tinimbang~$\pi$.
Hipotesis induksi maringi !!a{proof}~$\pi_1''$ kanggo
$!B, \Gamma \fCenter \Delta$. Saiki gatekna !!{proof}
kang pungkasané nganggo \CutCS{},
\[
\AxiomC{}
\RightLabel{$\pi_2'$}
\Deduce$\Gamma \fCenter \Delta, !B$
\AxiomC{}
\RightLabel{$\pi_1''$}
\Deduce$!B, \Gamma \fCenter \Delta$
\RightLabel{\CutCS}
\BinaryInf$\Gamma \fCenter \Delta$
\DisplayProof
\]
Pangkat cuté uga luwih cilik tinimbang pangkat cut~$\pi$,
mula hipotesis induksi maringi \Log{G3c}-!!{proof}~$\pi'$
kanggo $\Gamma \fCenter \Delta$.
\item Yèn $!A \ident \lforall[x][!B(x)]$, banjur $\pi$
pungkasané kaya mangkéné:
\[
\AxiomC{}
\RightLabel{$\pi_1'$}
\Deduce$\Gamma \fCenter \Delta, !B(c)$
\RightLabel{\RightR{\lforall}}
\UnaryInf$\Gamma \fCenter \Delta, \lforall[x][!B(x)]$
\LeftSubproofLabel{$\pi_1$}
\AxiomC{}
\RightLabel{$\pi_2'$}
\Deduce$!B(t), \lforall[x][!B(x)], \Gamma \fCenter \Delta$
\RightLabel{\LeftR{\lforall}}
\UnaryInf$\lforall[x][!B(x)], \Gamma \fCenter \Delta$
\RightSubproofLabel{$\pi_2$}
\RightLabel{\CutCS}
\BinaryInf$\Gamma \fCenter \Delta$
\DisplayProof
\]
Miturut hipotesis induksi, \CutCS{} bisa diilangi saka
\[
\AxiomC{}
\RightLabel{$\pi_1''$}
\Deduce$!B(t), \Gamma \fCenter \Delta, \lforall[x][!B(x)]$
\AxiomC{}
\RightLabel{$\pi_2'$}
\Deduce$!B(t), \lforall[x][!B(x)], \Gamma \fCenter \Delta$
\RightLabel{\CutCS}
\BinaryInf$!B(t), \Gamma \fCenter \Delta$
\DisplayProof
\]
ing ngendi $\pi_1''$ diolehké saka $\pi_1$ (dudu~$\pi_1'$!)
kanthi pelemahan kang njaga !!{height}. (!!{proof} iki
nduwèni pangkat cut kang padha nanging !!{height} cut
luwih cilik.) Mula kita ngolèhké !!a{proof}~$\pi_2''$
kanggo $!B(t), \Gamma \fCenter \Delta$.

!!{proof}~$\pi_1'$ nduwèni sekèn pungkasan
$\Gamma \Sequent \Delta, !B(c)$, ing ngendi $c$~variabel
eigen saka inferènsi~$\RightR{\lforall}$. Mula $c$~ora ana
ing~$!B(x)$, $\Gamma$, utawa~$\Delta$. Kita anggep
!!{proof}~$\pi$ reguler, mula $c$~mung dadi variabel eigen
ing inferènsi~$\RightR{\lforall}$ kang nyusul iki lan
mung ana ing sadhuwuré, yaiku mung ana ing~$\pi_1'$,
lan dudu variabel eigen sawijining inferènsi ing~$\pi_1'$.
Amarga $c$ ora ana ing~$\pi_2$, variabel iku ora bisa ana ing~$t$.
Miturut \olref[seq][qua]{lem:subst-regular},
!!{proof}~$\Subst{\pi_1'}{t}{c}$ iku !!a{proof} kanggo
$\Gamma \Sequent \Delta, !B(t)$. Saiki hipotesis induksi
ditrapaké marang !!{proof}
\[
\AxiomC{}
\RightLabel{$\Subst{\pi_1'}{t}{c}$}
\Deduce$\Gamma \fCenter \Delta, !B(t)$
\AxiomC{}
\RightLabel{$\pi_2''$}
\Deduce$!B(t), \Gamma \fCenter \Delta$
\RightLabel{\CutCS}
\BinaryInf$\Gamma \fCenter \Delta$
\DisplayProof
\]
(Hipotesis induksi bisa ditrapaké amarga !!{formula} cuté
yaiku~$!B(t)$, mula !!{proof} iku nduwèni pangkat cut
luwih cilik tinimbang~$\pi$.)
\item Kahanan $!A \ident \lexists[x][!B(x)]$ dipasrahaké
dadi latihan.
\end{enumerate}
\end{proof}

\begin{prob}
Priksanen subkahanan saka kahanan~D ing bukti
\olref{lem:cut-adm-G3c} nalika $!A$ dudu formula utama
ing inferènsi pungkasan~$\pi_1$ kang nganggo~$\LeftR{\lif}$.
(Cathetan: pérangan saka latihan iki yaiku njlentrehaké
apa kang kudu dibuktèkaké, yaiku wujudé $\pi$ ing kahanan iki.)
\end{prob}

\begin{prob}
Priksanen subkahanan saka kahanan~D ing bukti
\olref{lem:cut-adm-G3c} nalika $!A$ dudu formula utama
ing inferènsi pungkasan~$\pi_1$ kang nganggo
~$\LeftR{\lforall}$.
\end{prob}

\begin{prob}
Priksanen subkahanan saka kahanan~E ing bukti
\olref{lem:cut-adm-G3c} nalika $!A$ dadi formula utama
ing inferènsi pungkasan $\pi_1$ lan $\pi_2$,
lan $!A \ident (!B \land !C)$.
\end{prob}

\begin{prob}
Priksanen subkahanan saka kahanan~E ing bukti
\olref{lem:cut-adm-G3c} nalika $!A$ dadi formula utama
ing inferènsi pungkasan $\pi_1$ lan $\pi_2$,
lan $!A \ident \lexists[x][!B(x)]$.
\end{prob}

Gatekna yèn ing kahanan~E, !!{height} cut saka
!!{proof} kang pungkasané nganggo \CutCS{} lan kena
hipotesis induksi bisa luwih gedhé tinimbang !!{height}
cut saka bukti asal. Tuladhané, ing kahanan
$!A \ident (!B \lif !C)$, $\pi$ kita owahi dadi
!!a{proof} kang ngemot rong inferènsi \CutCS{},
\[
\AxiomC{}
\RightLabel{$\pi_2'$}
\Deduce$\Gamma \fCenter \Delta, !B$
\AxiomC{}
\RightLabel{$\pi_1'$}
\Deduce$!B, \Gamma \fCenter \Delta, !C$
\AxiomC{}
\RightLabel{$\pi_2''$}
\Deduce$!C, \Gamma \fCenter \Delta$
\RightLabel{\CutCS}
\BinaryInf$!B, \Gamma \fCenter \Delta$
\RightSubproofLabel{$\pi_1''$}
\RightLabel{\CutCS}
\BinaryInf$\Gamma \fCenter \Delta$
\DisplayProof
\]
\CutCS{} sisih ndhuwur ing antarané $\pi_1'$ lan $\pi_2''$
nduwèni pangkat cut lan !!{height} cut kang luwih cilik
tinimbang~$\pi$. Nanging !!{proof}~$\pi_1''$---asil saka
ngetrapaké hipotesis induksi---saiki ora mesthi nduwèni
!!{height} cut luwih cilik tinimbang~$\pi$. Amarga
!!{formula} cut saka \CutCS{} sisih ngisor yaiku $!B$,
\emph{pangkat} cuté tetep luwih cilik tinimbang pangkat
cut~$\pi$.

\olref{lem:cut-adm-G3c} nuduhaké yèn siji inferènsi
\CutCS{} kang paling dhuwur bisa diilangi saka !!a{proof}.
Kita bisa nuduhaké yèn pirang-pirang inferènsi \CutCS{}
ing !!a{proof} bisa diilangi kanthi ngetrapaké lema iku
kanthi urut marang sub-!!{proof}s paling dhuwur kang
pungkasané nganggo~\CutCS.

\begin{thm}\ollabel{thm:cut-elim-G3c-topmost} Saben !!{proof} $\pi$
ing $\Log{G3c} + \CutCS$ bisa diowahi dadi !!a{proof} ing~\Log{G3c}.
\end{thm}

\begin{proof}
  Buktiné nganggo induksi tumrap cacah~$n$ inferènsi
  \CutCS{} ing~$\pi$. Yèn $n=0$, ora ana kang perlu
  ditindakaké: $\pi$ wis dadi !!a{proof} ing~\Log{G3c}.
  Yèn ora, gatekna sub-!!{proof}~$\pi'$ kang pungkasané
  nganggo inferènsi \CutCS{} paling dhuwur. Ing kono
  mung ana siji aturan \CutCS{} minangka inferènsi
  pungkasan. Trapna \olref{lem:cut-adm-G3c} kanggo
  ngowahi dadi bukti tanpa \CutCS{}~$\pi''$, banjur
  gantenana sub-!!{proof}~$\pi'$ nganggo~$\pi''$.
  Asilé !!a{proof} kang ngemot $n-1$ inferènsi \CutCS{},
  mula hipotesis induksi bisa ditrapaké.
\end{proof}

\end{document}
